\section{Definitions, theorem and proof dependencies}\label{sec:main}
An alternating sign matrix has entries in $\{-1,0,1\}$, alternating
nonzero entries along each row and column, and all row and column sums
one. It is diagonally symmetric when it equals its transpose. Write
\[
 a_n=|\DS(n)|,\qquad S(A)=|\{i:A_{ii}\ne0\}|,\qquad
 Z_n(s)=\sum_{A\in\DS(n)}s^{S(A)},\qquad Z_0(s)=1.
\]
Thus $a_n=Z_n(1)$ and $a_n$ is A005163. All logarithms of positive
numbers are natural. Put
\begin{equation}\label{constants}
 c=\sqrt3,\quad \alpha=\tfrac12\log\frac{3\sqrt3}{4},\quad
 \beta=\tfrac14\log3,\quad \kappa=\tfrac5{72},\quad
 b=\tfrac34\log3-\tfrac12\log2.
\end{equation}
On $\HH_+=\{\Re s>0\}$ define
\begin{equation}\label{pressure}
 P(s)=\frac{(s+1)^2}{2(s+2)},\qquad
 B(s)=b+\tfrac12\Log P(s).
\end{equation}
Here $\Log P(s)=2\Log(s+1)-\Log(s+2)-\log2$, with the indicated
right-half-plane logarithms. In particular $P(c)=1$ and $B(1)=\beta$.
The companion polynomial and its centered logarithm are
\begin{equation}\label{companiondefs}
 D_n(t)=\frac{Z_n(\sqrt t)Z_{n+1}(\sqrt t)}{\sqrt t},\quad
 f(t)=\Log P(\sqrt t),\quad
 E_n(t)=\Log\frac{D_n(t)}{D_n(3)}-nf(t),
\end{equation}
where $t\in\Omega_t=\C\setminus(-\infty,0]$ and the square root
has positive real part. The quotient defining $D_n$ is interpreted
polynomially. Its logarithm is normalized to vanish at $t=3$.

\begin{theorem}[Full leading equivalent]\label{thm:main}
There is a holomorphic function $E$ on $\Omega_t$, real on $(0,\infty)$,
with $E(3)=0$, such that $E_n\to E$ locally uniformly. Define
\begin{equation}\label{Ccal}
 C_{\rm cal}=
 \frac{3^{11/72}\pi^{1/6}e^{\zeta'(-1)/6}}
      {2^{1/24}\Gamma(1/3)^{1/3}}>0
\end{equation}
and, on $\HH_+$,
\begin{equation}\label{amplitude}
 C(s)=C_{\rm cal}\exp\{E(s^2)/2\}
       (s/c)^{1/2}P(s)^{-1/4}.
\end{equation}
The powers use the analytic logarithms normalized at $s=c$.
Then $C$ is nowhere zero and positive for $s>0$, and
\begin{equation}\label{fullequivalent}
 Z_n(s)=C(s)e^{\alpha n^2+nB(s)}n^{-\kappa}(1+o(1))
\end{equation}
locally uniformly on $\HH_+$, along the full integer sequence.
In particular,
\begin{equation}\label{aequiv}
 a_n\sim C_{\rm cal}2^{-1/4}e^{E(1)/2}
               e^{\alpha n^2+\beta n}n^{-5/72}.
\end{equation}
A convergent operator/Fourier expression for $E(t)$ is given in
\eqref{Eexplicit}; it defines the constant without a numerical fit.
\end{theorem}

\subsection{What is supplied in the reader bundle}
The unchanged earlier article \cite{r124}, supplied both as a PDF and
as its complete TeX source in the bundle's \texttt{earlier\_input}
folder, proves the bounded-remainder stage. Its scope statements are
historical: its failure to assert a limit is strengthened here. All
new convergence, calibration and boundary steps are proved below;
no unpublished proof note is a required reader input.

For precision, these are the earlier results used, with their locations
in that supplied article:
\begin{enumerate}[label=\textup{(I\arabic*)}]
\item Sections 2--4: imaginary-axis zeros of $Z_n$, negative-real zeros
of $D_n$, cross-size interlacing, the exact adjacent determinant and
its finite Pascal reduction, and
$Z_n(c)Z_{n+1}(c)=c2^n A_{n+1}$, where $A_N=|\ASM(N)|$.
\item Sections 6--8: the exact degree-space compression
$D_n(t)/D_n(3)=\det(I+(t-3)P_n\mathcal B P_n)$, the bounded operator
$\mathcal B$, its unweighted decomposition $C_0+\mathcal E$ with
$0\le C_0\le I/3$, $\mathcal E\in\Stwo$,
$\mathcal S=(\mathcal E+\mathcal E^*)/2\in\Sone$, and the locally
uniform nonzero relative determinant $\mathcal A(t)$.
The exact continuous-Hahn/Fourier/Jacobi transform and the gamma-phase
comparison are included there.
\item Sections 9--11: the normalized Jacobi expansion, its endpoint
and derivative bounds, the mean formula
$\tr(J_nFJ_n)=(n+1)\pi^{-1}\int_0^\pi F+o(1)$ for endpoint-vanishing
$F$ with $\int |F|/\min(\theta,\pi-\theta)<\infty$, the global
angular kernel bound $|K_n(\theta,\phi)|\le C/|\theta-\phi|$, the
finite tilted-projection variance bound, and the Hilbert--Schmidt
weighted Hardy wave comparisons.
\item Sections 12 and 14: the convergent trace-class cross factor
$\Xi(t)$, the compact-uniform scalar bound
$\log T_n(t-3)-2nf(t)=O_K(1)$, the analogous companion bound, and
local uniform bounds on adjacent individual logarithmic differences.
\end{enumerate}
These are theorems and exact identities proved in \cite{r124}, rather
than assertions supplied by finite checks. We restate the definitions
needed to apply them. The additional published ingredients are the
ASM/OSASM product asymptotics in \cite{bfk}, the continuous-Hahn
normalization and generating function \cite{cll,dlmf}, and the
\emph{polynomial} Jacobi-statistic central limit theorem \cite{bd}.
We state precisely where each is used and prove the necessary extensions.

\subsection{The two logically separate limits}
An adjacent-product limit does not alone determine the individual
amplitudes: a residual $L+(-1)^n d$ has a constant adjacent sum.
Section~\ref{sec:companion} first proves $E_n\to E$.
Sections~\ref{sec:boundary}--\ref{sec:edge} then prove the distinct
boundary-ratio limit
\begin{equation}\label{Ftarget}
 F_n(s):=\frac{sZ_n(s)/Z_{n+1}(s)}{cZ_n(c)/Z_{n+1}(c)}
 \longrightarrow \frac{s/c}{\sqrt{P(s)}}.
\end{equation}
The finite-section proof is at the $1/n$ coefficient scale. In
particular, it does not substitute an infinite resolvent for a finite
one merely from strong convergence. Section~\ref{sec:consequences}
combines the two limits, proves the inverse threshold and derives
constant-order cumulants.
